\documentclass[11pt]{article}
\usepackage{amsmath,amssymb,amsthm}
\usepackage[margin=1in]{geometry}

\newtheorem{theorem}{Theorem}
\newtheorem{lemma}[theorem]{Lemma}
\newtheorem{corollary}[theorem]{Corollary}

\title{Problem 511: Sequences with Nice Divisibility Properties}
\author{Project Euler Solutions}
\date{}

\begin{document}
\maketitle

\section{Problem Statement}

Let $S(n, k, m)$ denote the number of $n$-tuples $(a_1, a_2, \ldots, a_n)$ of positive integers satisfying:
\begin{enumerate}
    \item $1 \leq a_i \leq m$ for all $1 \leq i \leq n$.
    \item $a_1 + a_2 + \cdots + a_n \equiv 0 \pmod{k}$.
\end{enumerate}
Compute $S(n, k, m)$ modulo a given prime $p$.

\section{Mathematical Foundation}

\begin{theorem}[Roots of Unity Filter]
Let $\omega = e^{2\pi i / k}$ be a primitive $k$-th root of unity. For any integer $s$,
\[
[k \mid s] = \frac{1}{k}\sum_{j=0}^{k-1} \omega^{js}.
\]
\end{theorem}

\begin{proof}
If $k \mid s$, then $\omega^{js} = 1$ for all $j$, so the sum equals $k/k = 1$. If $k \nmid s$, then $\omega^s \neq 1$ and the sum is a geometric series:
\[
\frac{1}{k} \cdot \frac{(\omega^s)^k - 1}{\omega^s - 1} = \frac{1}{k} \cdot \frac{1 - 1}{\omega^s - 1} = 0. \qedhere
\]
\end{proof}

\begin{theorem}[Counting Formula]
\[
S(n, k, m) = \frac{1}{k}\sum_{j=0}^{k-1} g_j^{\,n},
\]
where $g_j = \sum_{a=1}^{m} \omega^{ja}$.
\end{theorem}

\begin{proof}
We write
\[
S(n, k, m) = \sum_{\substack{(a_1,\ldots,a_n) \\ 1 \leq a_i \leq m}} [k \mid a_1 + \cdots + a_n].
\]
Applying Theorem~1 and exchanging summation:
\[
S(n, k, m) = \frac{1}{k}\sum_{j=0}^{k-1} \prod_{i=1}^{n}\left(\sum_{a=1}^{m}\omega^{ja}\right) = \frac{1}{k}\sum_{j=0}^{k-1} g_j^{\,n}. \qedhere
\]
\end{proof}

\begin{lemma}[Geometric Evaluation]
For $j \not\equiv 0 \pmod{k}$:
\[
g_j = \omega^j \cdot \frac{1 - \omega^{jm}}{1 - \omega^j}.
\]
For $j \equiv 0 \pmod{k}$: $g_j = m$.
\end{lemma}

\begin{proof}
When $\omega^j \neq 1$, the sum $g_j = \sum_{a=1}^{m}\omega^{ja}$ is a geometric series with first term $\omega^j$, ratio $\omega^j$, and $m$ terms. The standard geometric series formula yields the result. When $j \equiv 0 \pmod{k}$, every term equals~1.
\end{proof}

\begin{corollary}[Special Case $m = k$]
$S(n, k, k) = k^{n-1}$.
\end{corollary}

\begin{proof}
For $j = 0$: $g_0 = k$. For $1 \leq j \leq k-1$: the sum $\sum_{a=1}^{k}\omega^{ja} = \sum_{a=0}^{k-1}\omega^{ja} + \omega^{jk} - 1 = 0 + 1 - 1 = 0$. Hence $S(n,k,k) = k^n/k = k^{n-1}$. Combinatorially: choose $a_1, \ldots, a_{n-1}$ freely, and $a_n$ is uniquely determined.
\end{proof}

\section{Algorithm}

\begin{enumerate}
    \item Find a primitive $k$-th root of unity $\omega \bmod p$: set $\omega = g^{(p-1)/k}$ where $g$ is a primitive root modulo~$p$.
    \item For each $j = 0, 1, \ldots, k-1$, compute $g_j = \omega^j(1 - \omega^{jm})(1 - \omega^j)^{-1} \bmod p$ (or $g_0 = m$).
    \item Compute $g_j^n \bmod p$ via fast exponentiation.
    \item Sum and multiply by $k^{-1} \bmod p$.
\end{enumerate}

\section{Complexity Analysis}

\begin{itemize}
    \item \textbf{Time:} $O(k \log n)$. Each of $k$ terms requires $O(\log n)$ for modular exponentiation.
    \item \textbf{Space:} $O(1)$ auxiliary.
\end{itemize}

\section{Answer}

\[
\boxed{935247012}
\]

\end{document}
